
\documentclass[12pt, a4paper]{article}
\usepackage[utf8]{inputenc}
\usepackage[T2A]{fontenc}
\usepackage[left=2cm, right=2cm, top=2cm, bottom=2cm]{geometry}

\begin{document}

Непрерывность. Это понятие формализует идею, если $x_1$ и $x_2$ близки, то и $f(x_1)$ и $f(x_2)$ близки. Формально, для любого $\epsilon > 0$ существует $\delta > 0$, такое что если расстояние от $x_1$ до $x_2$ меньше, чем $\delta$, будет следовать то что расстояние от $f(x_1)$ до $f(x_2)$ меньше чем $\epsilon$. В произвольном пространстве близость это характеристика пары точек. В линейном пространстве есть операция вычитания. Это позволяет перенести вопрос о близости двух точек $x_1$ и $x_2$ к вопросу о близости к одному… к нулевому элементу. Таким образом, чтобы проверить непрерывность линейного оператора во всех точках, достаточно проверить его поведение в одной точке, в нуле.
\end{document}
